% Preparatory lemma: nilpotent-shift and binomial expansion
% Source: Wolf (2012), Chapter 8, lines 1211--1212
%
% This is the purely algebraic half: the closed formula for powers of a
% Jordan block.  The two-sided estimate on the largest singular value,
% Wolf Eq. (8.104), follows it below.
\begin{lemma}[Nilpotent shift and Jordan-block binomial expansion]
    \label{lem:jordan_block_power}
    \lean{Matrix.nilpotentShift_pow_apply, Matrix.nilpotentShift_pow_D_eq_zero, Matrix.jordanBlock_pow}
    \leanok
    This is the Jordan-block calculation in
    \cite[Chapter~8, proof of Equation~(8.104)]{Wolf2012Quantum}.
    Let $D\ge 1$, let $N\in M_{D}(\mathbb C)$ be the strict upper-shift matrix
    $N_{i,j}=1$ iff $j=i+1$ (and $0$ otherwise), and
    $J_{D}(\lambda)=\lambda I+N$ the associated Jordan block.
    Then $N^{k}_{i,j}=1$ iff $j=i+k$ (with $N^{k}=0$ for $k\ge D$), and for
    every $n\in\mathbb N$,
    \begin{align}
    J_{D}(\lambda)^{n}
      =\sum_{k=0}^{\min\{n,D-1\}}
        \binom{n}{k}\,\lambda^{\,n-k}\,N^{k}.
    \end{align}
\end{lemma}

\begin{proof}
  \leanok
  Because $\lambda I$ and $N$ commute, the ordinary binomial theorem expands
  $(\lambda I+N)^{n}$ into a sum of terms $\binom{n}{k}(\lambda I)^{n-k}N^{k}$
  for $0\le k\le n$.  Substituting $(\lambda I)^{n-k}=\lambda^{n-k}I$ and
  observing that $N^{k}=0$ whenever $k\ge D$ truncates the sum at
  $k\le\min\{n,D-1\}$.  The superdiagonal formula $N^{k}_{i,j}=1$ iff $j=i+k$
  is a straightforward induction on $k$ using the definition of $N$.
\end{proof}

% Source: Wolf (2012), Chapter 8, Equation (8.104), lines 1197--1215
\begin{lemma}[Two-sided estimate for a Jordan-block power]
    \label{lem:jordan_block_power_norm}
    \lean{Matrix.l2_opNorm_jordanBlock_pow_bounds}
    \leanok
    \uses{lem:jordan_block_power}
    This is \cite[Chapter~8, Equation~(8.104)]{Wolf2012Quantum}.
    Let $D\ge 1$, let $J_{D}(\lambda)=\lambda I+N$ be the Jordan block of
    Lemma~\ref{lem:jordan_block_power}, and write $\|\cdot\|_{\infty}$ for the
    largest singular value.  Then, for every $n\in\mathbb N$ and every
    $k_{0}\le\min\{n,D-1\}$,
    \begin{align}
      |\lambda|^{\,n-k_{0}}\binom{n}{k_{0}}
      \le\bigl\|J_{D}(\lambda)^{n}\bigr\|_{\infty}
      \le\sum_{k=0}^{\min\{n,D-1\}}|\lambda|^{\,n-k}\binom{n}{k}.
    \end{align}
\end{lemma}

\begin{proof}
    \leanok
    \uses{lem:jordan_block_power}
    By Lemma~\ref{lem:jordan_block_power} the power $J_{D}(\lambda)^{n}$ is the
    Toeplitz matrix $\sum_{k}\binom{n}{k}\lambda^{n-k}N^{k}$, whose first row
    carries the entries $\lambda^{n-k}\binom{n}{k}$.  Since the largest
    singular value dominates the modulus of every entry, reading the entry in
    row $0$ and column $k_{0}$ gives the left inequality; the hypothesis
    $k_{0}\le\min\{n,D-1\}$ is exactly what places that entry inside the
    matrix and inside the range of summation.  For the right inequality, the
    triangle inequality distributes the norm over the same expansion, and
    $\|N^{k}\|_{\infty}\le 1$ for every $k$, because $N^{\dagger}N$ is the
    diagonal projection $\operatorname{diag}(0,1,\dots,1)$ and the largest
    singular value is multiplicative under the adjoint pairing.
\end{proof}
